\documentclass[conference]{IEEEtran}

\title{a11y-TDD: An Accessibility-First Approach to Test-Driven Web Development}

\author{
  \IEEEauthorblockN{Josbel Luna}
  \IEEEauthorblockA{Department of Electrical and Computer Engineering\\
  University of Alberta}
  \and
  \IEEEauthorblockN{Kaparthini Reshmi}
  \IEEEauthorblockA{Department of Electrical and Computer Engineering\\
  University of Alberta}
  \and
  \IEEEauthorblockN{Onesh Punchinilame}
  \IEEEauthorblockA{Department of Electrical and Computer Engineering\\
  University of Alberta}
}

\begin{document}

\maketitle

\section{Background}

% Problem statement

Web accessibility is a core principle of web development that ensures websites can be used by everyone. 
This is important because the World Health Organization estimates that 1.3 billion people worldwide experience a significant disability~\cite{who2023disability}. 
Despite the relevance of this, web accessibility remains an unsolved problem.
Martins and Duarte found that only 0.5\% of 2.88 million web pages had no automatically detectable accessibility errors~\cite{martins2024accessibility}.

% LLM generated code for web
With the rise of LLMs, one might hope that they could solve web accessibility issues by following those rules closely, but this is far from reality.
In a study where developers prompted LLMs to generate websites, Aljedaani et al. reported that 84\% of the resulting websites had accessibility issues~\cite{aljedaani2024chatgpt}.
Related research by Guriță and Vatavu also demonstrated that although accessibility-oriented prompts reduced violations, semantic accessibility barriers remained~\cite{gurita2025barriers}.
 
% LLM for TDD



% [Research x] says that LLM generated tests for TDD are good fo code.

% These studies examine either accessibility-oriented prompting or test-driven code generation, but there appears to be limited work evaluating whether accessibility tests used as explicit TDD requirements cause LLMs to generate more accessible web applications than ordinary prompting.

% Ideas:


\section{Planned Approach}

\section{Planned Evaluation}

\section{Project Timeline}

\section*{References}

\begin{thebibliography}{1}
  \bibitem{who2023disability}
World Health Organization, ``Disability,'' Mar. 7, 2023. [Online]. Available: https://www.who.int/news-room/fact-sheets/detail/disability-and-health

  \bibitem{martins2024accessibility}
  B. Martins and C. Duarte, ``A large-scale web accessibility analysis considering technology adoption,'' \emph{Universal Access in the Information Society}, vol. 23, no. 4, pp. 1857--1872, 2024, doi: 10.1007/s10209-023-01010-0.

  \bibitem{aljedaani2024chatgpt}
  W. Aljedaani, A. Habib, A. Aljohani, M. M. Eler, and Y. Feng, ``Does ChatGPT generate accessible code? Investigating accessibility challenges in LLM-generated source code,'' in \emph{Proc. 21st Int. Web for All Conf. (W4A '24)}, 2024, pp. 165--176, doi: 10.1145/3677846.3677854.

  \bibitem{gurita2025barriers}
  A.-E. Guriță and R.-D. Vatavu, ``When LLM-generated code perpetuates user interface accessibility barriers, how can we break the cycle?'' in \emph{Proc. 22nd Int. Web for All Conf. (W4A '25)}, 2025, pp. 124--134, doi: 10.1145/3744257.3744266.
\end{thebibliography}


\end{document}
